%%% --- Общие настройки документа ---
\usepackage{geometry}
\geometry{left=20mm,right=20mm,top=20mm,bottom=20mm} % поля
\usepackage{listings}
\lstset{
  basicstyle=\ttfamily\footnotesize,
  breaklines=true,
  breakatwhitespace=false,
  postbreak=\mbox{\textcolor{red}{$\hookrightarrow$}\space},
  frame=single,
  numbers=left,
  numberstyle=\tiny,
  numbersep=5pt,
  tabsize=2,
  showstringspaces=false,
  keywordstyle=\color{blue},
  commentstyle=\color{gray},
  stringstyle=\color{orange},
  columns=flexible
}
\usepackage{cmap}                     % поиск в PDF
\usepackage[T2A]{fontenc}             % кодировка шрифтов
\usepackage[utf8]{inputenc}           % кодировка исходного текста
\usepackage[english, russian]{babel}  % локализация и переносы

\usepackage{xcolor}
\usepackage{hyperref}
\hypersetup{
    pdfstartview=FitH,
    colorlinks=true,
    linkcolor=blue,
    filecolor=magenta,
    urlcolor=cyan,
}

%%% --- Математика ---
\usepackage{amsmath,amssymb,amsthm,mathtools,bm}
\usepackage{icomma}                   % умная запятая в числах
\allowdisplaybreaks[4]                % перенос формул
\binoppenalty=1000                    % уменьшение переносов в формулах

%%% --- Символы и шрифты ---
\usepackage{euscript}    % Евклид
\usepackage{mathrsfs}    % Красивый матшрифт

%%% --- Интервалы и абзацы ---
% \sloppy                              % строго соблюдать границы текста (отключено)
\setlength{\emergencystretch}{3em}   % предотвращает overfull hbox
\linespread{1.2}                     % межстрочный интервал
\setlength{\parskip}{0pt}            % интервал между абзацами
\setlength{\parindent}{0pt}          % без красной строки

%%% --- Работа с рисунками ---
\usepackage{graphicx}
\graphicspath{{images/}}
\usepackage{wrapfig}
\usepackage{caption}
\captionsetup{justification=centering} % подписи по центру

%%% --- Работа с таблицами ---
\usepackage{array,tabularx,tabulary,booktabs,longtable,multirow}

%%% --- Списки ---
\usepackage{enumitem}
\makeatletter
\AddEnumerateCounter{\asbuk}{\russian@alph}{щ}
\makeatother

%%% --- TikZ и PGFPlots ---
\usepackage{tikz,pgfplots}
\usetikzlibrary{arrows,graphs,graphs.standard}
\pgfplotsset{width=10cm,compat=newest}
\pgfkeys{/pgf/trig format=rad}

%%% --- Прочее ---
\usepackage{float}    % [H] для фиксации положения
\usepackage{gensymb}  % красивый знак градусов
\usepackage[normalem]{ulem} % подчёркивания и зачёркивания

%%% --- Колонтитулы ---
\usepackage{titleps}
\newpagestyle{main}{
  \sethead{}{}{}
  \setfoot{}{\thepage}{}
}
\pagestyle{main}

%%% --- Оформление разделов ---
\setcounter{secnumdepth}{0} % отключаем нумерацию секций